\section{Problem Definition}
\label{sec:problem}
% Input object, decision, policy, optimization problem. Each definition is followed by the running example.

\begin{definition}[Candidate Set]
\label{def:candidate}
Given a table $T$, the candidate set $\cand$ is <definition>.
\end{definition}

\begin{example}
\label{ex:candidate}
In Example~\ref{ex:motivation}, <the candidate set of the running example>.
\end{example}

\begin{definition}[Labeling Policy]
\label{def:policy}
A policy $\policy$ maps <state> to <decision> such that at most $\budget$ candidates are labeled.
\end{definition}

The goal is
\begin{equation}
\label{eq:objective}
\max_{\policy}\ \mathrm{Quality}(\policy(\cand)) \quad \text{s.t.} \quad |\policy(\cand)| \le \budget .
\end{equation}

\begin{theorem}[Hardness]
\label{thm:hard}
Under general per-candidate costs, the decision version of Problem~\eqref{eq:objective} is NP-complete.
\end{theorem}
Appendix~\ref{app:proof-hard} gives the proof. The result motivates the approximation of Section~\ref{sec:method}.
